% Experimental setup (Task 42): old Setup without the constraint model (now in Section 3) and without the two
% limitations tables (now Appendix C); T1 and T2 merged into t_data.
\section{Experimental Setup}
\label{sec:setup}

\paragraph{Datasets.}
We used LYCOS-IDS2017 and LycoS-Unicas-IDS2018 (called 2017 and 2018 below), built from the traffic of CIC-IDS2017
and CSE-CIC-IDS2018 \cite{sharafaldin2018,aws_cse_cic_ids2018}, both with the same \nFeatures{} flow features (Table~\ref{t_data}).
On 2017 we removed Heartbleed (\nHeartbleed{} rows) and the structured query language (SQL) injection web attack (\nSqlInjection{} rows, all after
the temporal cut), merged the remaining web attacks, removed feature vectors that carried more than one label and
exact duplicates, and split each (label, day) group in time with gaps between train, validation and test.
On 2018 we removed FTP-Patator (\nFtpRemoved{} rows) and DoS Slowhttptest (\nSlowhttpRemoved{} rows) because the two
classes were indistinguishable in the features, merged the DDoS and web attack variants, removed conflicting and
duplicate vectors, capped large classes, and split at random with stratification, because the file has no timestamps.
This left \nClassesSeventeen{} classes on 2017 and \nClassesEighteen{} on 2018; all cross-dataset comparisons use the
\nCommonClasses{} common attack classes.

% Generated by paper_v3/gen/make_tables_v3.py on 2026-10-30. Do not edit.
% Source: results/stage1_class_counts.csv
% Source: results/stage1_cleaning_report.csv
% Source: results/stage7/stage7b_counts_by_step.csv
% Source: results/stage2_per_class_f1.csv
% Source: results/stage7/stage7b_per_class.csv
\begin{table}[t]
\centering
\footnotesize
\caption{Datasets after cleaning (rows per class and split) and test macro-F1 of the models.}
\label{t_data}
\sbox0{%
\begin{tabular}{lrrrrrr}
\toprule
 & \multicolumn{3}{c}{LYCOS-IDS2017} & \multicolumn{3}{c}{LycoS-Unicas-IDS2018} \\
Class & Train & Val & Test & Train & Val & Test \\
\midrule
Benign & 740{,}728 & 246{,}910 & 246{,}908 & 600{,}000 & 200{,}000 & 200{,}000 \\
Bot & 431 & 144 & 144 & 57{,}660 & 19{,}220 & 19{,}220 \\
DDoS & 56{,}261 & 18{,}754 & 18{,}754 & 180{,}000 & 60{,}000 & 60{,}000 \\
DoS GoldenEye & 3{,}977 & 1{,}326 & 1{,}326 & 16{,}116 & 5{,}372 & 5{,}373 \\
DoS Hulk & 93{,}476 & 31{,}159 & 31{,}158 & 180{,}000 & 60{,}000 & 60{,}000 \\
DoS Slowhttptest & 2{,}770 & 923 & 923 & removed & removed & removed \\
DoS Slowloris & 3{,}080 & 1{,}027 & 1{,}027 & 5{,}428 & 1{,}809 & 1{,}809 \\
FTP-Patator & 2{,}351 & 784 & 783 & removed & removed & removed \\
PortScan & 981 & 327 & 327 & absent & absent & absent \\
SSH-Patator & 1{,}739 & 580 & 580 & 55{,}571 & 18{,}524 & 18{,}524 \\
Web attack & 1{,}187 & 395 & 397 & 256 & 85 & 85 \\
\midrule
Total & 906{,}981 & 302{,}329 & 302{,}327 & 1{,}095{,}031 & 365{,}010 & 365{,}011 \\
\midrule
\multicolumn{7}{l}{\textit{Test macro-F1}} \\
FT-Transformer & \multicolumn{3}{c}{$0.9231$ {\scriptsize ($0.9133$ to $0.9309$)}} & \multicolumn{3}{c}{$0.9924$ {\scriptsize ($0.9915$ to $0.9936$)}} \\
XGBoost & \multicolumn{3}{c}{$0.9515$} & \multicolumn{3}{c}{$0.9977$} \\
\midrule
\multicolumn{7}{l}{\textit{Removed labels (rows before cleaning) and reason}} \\
\multicolumn{7}{l}{2017: Heartbleed ($11$ rows): too few rows for a train / val / test split} \\
\multicolumn{7}{l}{2017: Web Attack SQL Injection ($12$ rows): all rows fall after the temporal cut (test only)} \\
\multicolumn{7}{l}{2018: FTP-Patator ($190{,}300$ rows): indistinguishable from DoS Slowhttptest in the 72 features} \\
\multicolumn{7}{l}{2018: DoS Slowhttptest ($105{,}550$ rows): indistinguishable from FTP-Patator in the 72 features} \\
\bottomrule
\end{tabular}}%
\ifdim\wd0>\textwidth\resizebox{\textwidth}{!}{\usebox0}\else\usebox0\fi
\par\smallskip\parbox{\linewidth}{\scriptsize 2017: temporal split per (label, day) group with gaps between splits. 2018: stratified random split (the file has no timestamps), with per-class caps applied before the split. Removed or absent classes are not modelled on that dataset. Macro-F1: FT-Transformer mean over 3 seeds (seed range), XGBoost one model; same training recipe on both datasets, no tuning on 2018.}
\end{table}


\paragraph{Models.}
We trained an FT-Transformer with \nSeeds{} seeds and an extreme gradient boosting (XGBoost) reference with one recipe on both datasets, without
tuning on 2018.
Test macro-F1 of the FT-Transformer was \fOneFTTSeventeen{} (seed range \fOneFTTSeventeenRange) on 2017 and
\fOneFTTEighteen{} (\fOneFTTEighteenRange) on 2018; XGBoost reached \fOneXGBSeventeen{} and \fOneXGBEighteen{}
(Table~\ref{t_data}).

\paragraph{Explanation and attack sets.}
Clean explanations used up to \paramExplPerClass{} test rows per class that all seeds classified correctly.
Attacks and detectors used the first \paramPerClass{} of these rows per attack class (fewer when a class had fewer;
web attack on 2018 has \nWebTestEighteen{} test rows), the same rows for every seed.

\paragraph{Near-duplicate filter and primary row set.}
For every test row we computed the exact Euclidean distance to its nearest training row in the standardised feature
space of its dataset and flagged rows with a distance below \paramNDThr.
These flags covered \ndSlowlorisSeventeen{} of the DoS Slowloris and \ndWebSeventeen{} of the web attack test rows on
2017, \ndSlowlorisEighteen{} of the DoS Slowloris rows on 2018, and \ndAllSeventeen{} and \ndAllEighteen{} of all test
rows (Appendix~\ref{sec:sens}).
The 2018 attack source rows excluded flagged DoS Slowloris rows by design, while the 2017 attack and detector source
rows were drawn without this filter: on the \nCommonClasses{} common classes they contained \ndSrcCommonSeventeen{}
flagged rows per seed.
We report three row sets: V0 (all source rows), V1 (flagged source rows removed in both datasets) and V2
(DoS Slowloris and web attack removed in both datasets, \nVtwoClasses{} classes).
V1 is the primary set for all results; V0 and V2 are sensitivity analyses (Section~\ref{sec:rq3} and
Appendix~\ref{sec:sens}).
The filter was applied after the attacks and after detector training, by removing source rows from the saved results;
models, attacks and detector thresholds are unchanged.
V1 was fixed as the primary set after the sensitivity results were seen (Appendix~\ref{sec:sens}), which is why V0 and
V2 are reported as well.
Per seed, V1 contains \nSrcVoneSeventeen{} source rows on 2017 (V0: \nSrcVzeroSeventeen) and \nSrcVoneEighteen{} on
2018, where no source row was flagged.

\paragraph{Seeds and uncertainty.}
All reported analyses used seeds 0, 1 and 2 on both datasets (Appendix~\ref{sec:sens}); a preliminary
detector run with seed 0 alone, used during development, is not reported.
We report means with the range over seeds, and for paired detector comparisons \paramCILevel{} bootstrap intervals
(\paramBoot{} draws): within seed for the equal-effort union, and pooled over seeds for the paired true-positive-rate
(TPR) differences.
Neither interval captures variation between seeds, so we give the per-seed values next to every headline difference.
